% ---------------------------------------------------------------------------
% Shared preamble for the Deep-Dive Compendium (loaded by compendium.tex and
% by the standalone chapter test wrapper). Chapter files must NOT load packages.
% ---------------------------------------------------------------------------
\usepackage[T1]{fontenc}
\usepackage[utf8]{inputenc}
\usepackage{mathpazo}
\linespread{1.04}
\renewcommand{\sfdefault}{phv}
\renewcommand{\ttdefault}{pcr}
\usepackage[letterpaper,margin=1in,headheight=14pt]{geometry}
\usepackage{microtype}
\usepackage{amsmath,amssymb,amsthm,mathtools,bm}
\usepackage{booktabs,array,longtable,tabularx,multirow}
\usepackage{enumitem}
\usepackage{xcolor}
\usepackage{graphicx}
\usepackage[most]{tcolorbox}
\usepackage{fancyhdr}
\usepackage{titlesec}
\usepackage{url}
\usepackage[hidelinks,pdfusetitle,bookmarksnumbered]{hyperref}
\usepackage[capitalise,noabbrev]{cleveref}

\setlist{itemsep=2pt,topsep=3pt,parsep=0pt}
\setlength{\parskip}{3pt}
\setlength{\emergencystretch}{2em}
\urlstyle{same}
\allowdisplaybreaks

% ---- Chapter and section headings ----
\titleformat{\chapter}[display]
  {\normalfont\huge\bfseries}{\chaptertitlename\ \thechapter}{12pt}{\Huge}
\titlespacing*{\chapter}{0pt}{-10pt}{24pt}

% ---- Headers ----
\pagestyle{fancy}
\fancyhf{}
\fancyhead[L]{\small\nouppercase{\leftmark}}
\fancyhead[R]{\small Deep-Dive Compendium, compiled 2026-09-07}
\fancyfoot[C]{\thepage}
\renewcommand{\headrulewidth}{0.3pt}
\fancypagestyle{plain}{\fancyhf{}\fancyfoot[C]{\thepage}\renewcommand{\headrulewidth}{0pt}}

% ---- Theorem-like environments ----
\theoremstyle{plain}
\newtheorem{theorem}{Theorem}[chapter]
\newtheorem{proposition}[theorem]{Proposition}
\newtheorem{lemma}[theorem]{Lemma}
\newtheorem{corollary}[theorem]{Corollary}
\newtheorem{law}[theorem]{Law}
\newtheorem{conjecture}[theorem]{Conjecture}
\theoremstyle{definition}
\newtheorem{definition}[theorem]{Definition}
\newtheorem{assumption}[theorem]{Assumption}
\newtheorem{example}[theorem]{Example}
\newtheorem{construction}[theorem]{Construction}
\theoremstyle{remark}
\newtheorem{remark}[theorem]{Remark}

% ---- Colours ----
\definecolor{provgray}{RGB}{245,245,245}
\definecolor{keyblue}{RGB}{232,240,250}
\definecolor{keyblueframe}{RGB}{60,100,160}
\definecolor{revorange}{RGB}{255,244,229}
\definecolor{revorangeframe}{RGB}{200,110,20}
\definecolor{bgframe}{RGB}{120,120,120}

% ---- Boxes ----
% Provenance block at the top of each chapter
\newtcolorbox{provenance}{enhanced,breakable,colback=provgray,colframe=bgframe,
  boxrule=0.4pt,arc=1pt,left=6pt,right=6pt,top=4pt,bottom=4pt,
  fontupper=\small,title=\small\bfseries Provenance,coltitle=black,colbacktitle=provgray,
  attach boxed title to top left={yshift=-2mm,xshift=4mm},boxed title style={boxrule=0pt,colframe=provgray}}
% Key ideas summary box
\newtcolorbox{keyideas}{enhanced,breakable,colback=keyblue,colframe=keyblueframe,
  boxrule=0.6pt,arc=1pt,left=6pt,right=6pt,top=4pt,bottom=4pt,
  title=\bfseries Key ideas in this chapter,coltitle=white}
% Reviewer's note (compiler's own flags)
\newtcolorbox{reviewernote}[1][]{enhanced,breakable,colback=revorange,colframe=revorangeframe,
  boxrule=0.6pt,arc=1pt,left=6pt,right=6pt,top=4pt,bottom=4pt,fontupper=\small,
  title=\small\bfseries Reviewer's note #1,coltitle=white}

% ---- Status tags for results (mirror the dives' own labels) ----
\newcommand{\tagEstablished}{\textcolor{keyblueframe}{\textsf{\small[established in the dive]}}}
\newcommand{\tagTransported}{\textcolor{keyblueframe}{\textsf{\small[transported from the literature]}}}
\newcommand{\tagNumerical}{\textcolor{keyblueframe}{\textsf{\small[numerically established]}}}
\newcommand{\tagConjecture}{\textcolor{revorangeframe}{\textsf{\small[conjecture / speculative]}}}
\newcommand{\tagRefuted}{\textcolor{revorangeframe}{\textsf{\small[refuted within the dive]}}}

% ---- Cross-reference helpers ----
\newcommand{\dive}[1]{Dive~#1}
\newcommand{\BL}[1]{BL\,#1}

% ---- Math macros ----
\newcommand{\E}{\mathbb{E}}
\newcommand{\Prob}{\mathbb{P}}
\newcommand{\R}{\mathbb{R}}
\newcommand{\N}{\mathbb{N}}
\DeclareMathOperator{\Var}{Var}
\DeclareMathOperator{\Cov}{Cov}
\DeclareMathOperator{\Corr}{Corr}
\DeclareMathOperator{\tr}{tr}
\DeclareMathOperator{\diag}{diag}
\DeclareMathOperator{\rank}{rank}
\DeclareMathOperator{\sign}{sign}
\DeclareMathOperator*{\argmin}{arg\,min}
\DeclareMathOperator*{\argmax}{arg\,max}
\newcommand{\ind}{\mathbf{1}}
\newcommand{\one}{\mathbf{1}}
\newcommand{\T}{\mathsf{T}}
\newcommand{\abs}[1]{\left|#1\right|}
\newcommand{\norm}[1]{\left\|#1\right\|}
\newcommand{\inner}[2]{\left\langle #1,#2\right\rangle}
\newcommand{\dd}{\mathrm{d}}
\newcommand{\ROAS}{\mathrm{ROAS}}
\newcommand{\mROAS}{\mathrm{mROAS}}
\newcommand{\iROAS}{\mathrm{iROAS}}
